% Village Life Needs: A Comprehensive Framework for Nepal
% Context: From remote rural (no roads, no electricity) to fully connected villages
% Think large. Think everything.

\documentclass[11pt,a4paper]{article}
\usepackage[utf8]{inputenc}
\usepackage[T1]{fontenc}
\usepackage{geometry}
\usepackage{enumitem}
\usepackage{titlesec}
\usepackage{hyperref}
\usepackage{parskip}

\geometry{margin=2.5cm}
\setlist{nosep, leftmargin=*}

\title{Village Life Needs: A Comprehensive Framework for Nepal}
\author{Context: Rural Villages --- Remote to Fully Connected}
\date{}

\begin{document}

\maketitle
\tableofcontents
\newpage

%==============================================================================
\section{Introduction}
%==============================================================================

This document catalogs the essential needs for village life in Nepal, considering the vast spectrum of conditions across the country. Nepal's topography---from the high Himalayas to the Terai plains---and varying levels of development mean that villages range from \textbf{fully remote} (no motor roads, no electricity, often accessible only by foot or mule) to \textbf{fully connected} (roads, grid electricity, internet). Each category below considers both extremes and the continuum between them.

\textbf{Think large.} This framework aims to be exhaustive---covering not only survival but dignity, not only basics but aspiration, not only individuals but community, not only present but future.

This document draws on research articles, policy reports, and documented studies on rural Nepal and village development (see References and Further Reading).

%==============================================================================
\section{Water}
%==============================================================================

\subsection{Drinking \& Domestic}
\begin{itemize}
    \item Springs (\textit{kuwa}, \textit{dhara}), wells, rivers, streams
    \item Rainwater harvesting, rooftop collection
    \item Piped supply, gravity-fed systems
    \item Water purification: boiling, filters, chlorine, UV
    \item Storage: tanks, \textit{ghagri}, \textit{gagri}
    \item Traditional: \textit{ghaite} (stone spouts), \textit{hitri}
    \item Water user committees, maintenance
    \item Seasonal scarcity: dry season, monsoon turbidity
    \item Water quality testing
\end{itemize}

\subsection{Irrigation}
\begin{itemize}
    \item Canals (\textit{kulo}), diversion weirs
    \item Lift irrigation, pump irrigation
    \item Drip, sprinkler (where feasible)
    \item Traditional \textit{pani ghatta} (water mills)
    \item Pond, reservoir
    \item Snowmelt, glacier-fed (high altitude)
\end{itemize}

\subsection{Irrigation Equipment \& Machinery (Bare Minimum)}
To make irrigation easier---especially where gravity-fed canals are insufficient or land needs preparation---villages need at least the following.

\textbf{Water lifting \& conveyance:}
\begin{itemize}
    \item \textbf{Pump:} Diesel, electric, or solar---to lift water from well, river, stream, pond when source is below field
    \item \textbf{Pipes \& hose:} PVC, HDPE, flexible hose---to convey water from source to field; suction and delivery pipes
    \item \textbf{Hand/treadle pump:} Where no fuel or electricity---human-powered lift
\end{itemize}

\textbf{Land preparation \& earthwork:}
\begin{itemize}
    \item \textbf{Power tiller (2-wheel tractor):} Plowing, tilling, land leveling---levels field for even water distribution, reduces runoff; common in Nepal Terai and accessible hills
    \item \textbf{Tractor (4-wheel):} Where terrain and road allow---plowing, tilling, pulling trailer, sometimes pump; heavier work
    \item \textbf{Attachments:} Plow, cultivator, leveler, ridger---for tiller/tractor
    \item \textbf{Shovel, spade, \textit{kodalo}:} Manual canal digging, repair---always needed; bare minimum when no machinery
\end{itemize}

\textbf{Canal \& distribution:}
\begin{itemize}
    \item \textbf{Canal lining:} Cement, plastic sheet---reduces seepage, saves water
    \item \textbf{Gates, valves:} To control flow, allocate water
    \item \textbf{Drip kit:} Main line, laterals, emitters---efficient for vegetables, orchards; saves water
    \item \textbf{Sprinkler:} Pipes, sprinkler heads---for larger area; needs pressure
\end{itemize}

\textbf{Context:}
\begin{itemize}
    \item \textbf{No road:} Tractor/tiller cannot reach; manual tools, solar pump (portable), or gravity only
    \item \textbf{With road:} Power tiller, tractor, diesel pump---all feasible; shared or cooperative ownership common; FMIS (farmer-managed irrigation) accounts for ~2/3 of Nepal's irrigated area and often outperforms agency-managed systems (see References)
    \item \textbf{With electricity:} Electric pump, solar pump---no fuel dependency
    \item \textbf{Maintenance:} Spare parts, repair skill---pump and machinery need upkeep
\end{itemize}

\subsection{Emergency \& Backup}
\begin{itemize}
    \item Hand pumps, tube wells
    \item Community water points
    \item Emergency tankers (disaster)
\end{itemize}

%==============================================================================
\section{Fire, Fuel \& Energy}
%==============================================================================

\subsection{Cooking Fuel}
\begin{itemize}
    \item Firewood: primary in remote areas; sustainable harvesting, community forests
    \item Improved cookstoves (\textit{improved chulo}), smoke hoods
    \item Biogas: where livestock exists
    \item LPG/gas cylinders: where roads allow delivery
    \item Kerosene: backup; lanterns, stoves
    \item Solar cookers: experimental
    \item Electric stoves: where reliable electricity exists
    \item Agricultural residue: straw, dung cakes
\end{itemize}

\subsection{Electricity}
\begin{itemize}
    \item \textbf{Off-grid:} Solar home systems (SHS), micro-hydro, pico-hydro---~3,300 micro-hydro plants in rural Nepal; research ranks micro-hydro best for rural electrification, then solar (see References)
    \item Solar lanterns, solar charging stations
    \item Grid connection (often unreliable)
    \item Micro-grid, mini-grid
    \item Backup: generators (diesel), batteries
    \item Street lighting
    \item Community charging points
\end{itemize}

\subsection{Mechanical \& Productive Use}
\begin{itemize}
    \item Grinding mills (\textit{ghatta}), rice hullers
    \item Irrigation pumps
    \item Welding, carpentry power tools
    \item Refrigeration, cold storage
    \item Human \& animal power (where no electricity)
\end{itemize}

\subsection{Vehicular \& Machine Fuel}
\begin{itemize}
    \item Petrol, diesel (where roads exist)
    \item Kerosene for generators
    \item Electric vehicles, charging (emerging)
\end{itemize}

%==============================================================================
\section{Food}
%==============================================================================

\subsection{Agriculture}
\begin{itemize}
    \item Seeds: local, improved, hybrid; seed banks---research in mid-hills Nepal shows improved cultivars and crop yields as primary drivers of household food self-sufficiency (Karki et al., 2015; see References)
    \item Fertilizers: organic (compost, manure), chemical (urea, DAP)
    \item Pesticides, herbicides
    \item Tools: plough (\textit{halo}), sickle (\textit{hasiya}), \textit{kodalo}, \textit{khurpi}, spade
    \item Terraces, land leveling
    \item Agricultural extension, training
    \item Weather information, crop calendar
    \item Pest control, crop protection
\end{itemize}

\subsection{Staples}
\begin{itemize}
    \item Rice, wheat, maize, millet, barley, buckwheat
    \item Varies by altitude: Terai (rice), hills (maize, millet), high hills (buckwheat, potato)
    \item Pulses, legumes
\end{itemize}

\subsection{Vegetables, Fruits \& Other}
\begin{itemize}
    \item Kitchen gardens, seasonal produce
    \item Fruit trees: citrus, mango, banana, apple (high altitude)
    \item Spices, herbs
    \item Storage: drying, pickling, smoking
\end{itemize}

\subsection{Livestock}
\begin{itemize}
    \item Buffalo, cows, goats, sheep, chickens, pigs, ducks
    \item Milk, meat, eggs, manure, draft power
    \item Fodder, grazing land
    \item Veterinary services, vaccination
    \item Beekeeping, fish farming (ponds)
\end{itemize}

\subsection{Storage \& Processing}
\begin{itemize}
    \item Grain storage (\textit{bharako}), rodent protection
    \item Cold storage (where electricity exists)
    \item Rice mills, oil mills, flour mills
    \item Cheese, ghee, \textit{chyurpi} (dried cheese)
    \item Local processing vs. distant
\end{itemize}

\subsection{Markets \& Distribution}
\begin{itemize}
    \item Local shops, \textit{haat bazaar}, weekly markets
    \item Access to district headquarters, wholesale
    \item Fair price, exploitation
    \item Food preservation for off-season
\end{itemize}

\subsection{Nutrition \& Security}
\begin{itemize}
    \item Dietary diversity
    \item Fortified foods, iodized salt
    \item School meals, mid-day meal
    \item Emergency food, relief
    \item Community kitchens (disaster)
    \item Malnutrition programs
\end{itemize}

%==============================================================================
\section{Shelter \& Housing}
%==============================================================================

\subsection{Construction}
\begin{itemize}
    \item Materials: stone, mud, wood, bamboo, brick, tin roofing
    \item Earthquake-resistant construction
    \item Retrofitting of old houses
    \item Ventilation, chimney (smoke)
    \item Insulation (cold areas)
\end{itemize}

\subsection{Comfort}
\begin{itemize}
    \item Fireplace (\textit{chulo}), warmth
    \item Blankets, bedding
    \item Lighting: sunlight, fire, kerosene, solar, electric
\end{itemize}

\subsection{Facilities}
\begin{itemize}
    \item Toilet: pit latrine, pour-flush, septic
    \item Bathing area
    \item Kitchen, storage
\end{itemize}

\subsection{Resettlement}
\begin{itemize}
    \item Landslide, flood relocation
    \item Planned resettlement
    \item Temporary shelter (disaster)
\end{itemize}

%==============================================================================
\section{Clothing \& Textiles}
%==============================================================================

\begin{itemize}
    \item Fabric, cloth; local weaving vs. market
    \item Tailoring, sewing, repair
    \item Seasonal: winter, monsoon, summer
    \item Traditional wear: \textit{daura suruwal}, \textit{sari}, \textit{gunyu}
    \item Footwear: sandals, shoes, boots
    \item School uniforms
    \item Washing, drying
\end{itemize}

%==============================================================================
\section{Transportation}
%==============================================================================

\subsection{No Motor Roads}
\begin{itemize}
    \item Foot trails, \textit{ghat} (stone steps)---research finds road access is the highest infrastructure priority in rural Nepal; remote areas show 50--65\% lower per capita income (see References)
    \item Trail maintenance, landslide clearance
    \item Pack animals: mules, horses, yaks (high altitude)
    \item Porters: human carriers
    \item Ropeways: gravity goods (\textit{tatopani}), passenger
    \item Rivers: boats, rafts, canoe (Terai, hills)
    \item Bridges: suspension, wooden, rope
    \item Helicopter: emergency only
\end{itemize}

\subsection{With Motor Roads}
\begin{itemize}
    \item Buses, jeeps, trucks, tractors
    \item Motorcycles, bicycles
    \item Ambulance, emergency vehicles
    \item Goods delivery
    \item Bus stops, parking
\end{itemize}

\subsection{Infrastructure}
\begin{itemize}
    \item Road construction, maintenance
    \item Culverts, drainage
    \item Retaining walls
    \item Ferry services (rivers)
    \item STOL airstrips (remote)
\end{itemize}

%==============================================================================
\section{Education}
%==============================================================================

\subsection{Early Childhood}
\begin{itemize}
    \item Daycare, \textit{bal bikas} centers
    \item ECD (Early Childhood Development)
    \item Mothers' groups, community-run
\end{itemize}

\subsection{School}
\begin{itemize}
    \item Primary (1--5), secondary (6--10), higher secondary (11--12)
    \item Buildings, classrooms, toilets, drinking water
    \item Teacher availability, training, housing
    \item Multi-grade teaching (remote)
    \item Textbooks, school supplies, uniforms
    \item Mid-day meal
    \item Girls' education, menstrual hygiene
    \item School management committee
\end{itemize}

\subsection{Higher \& Vocational}
\begin{itemize}
    \item Colleges, universities
    \item Vocational: carpentry, tailoring, agriculture, plumbing, electrician
    \item Distance learning (internet)
    \item Scholarships, hostel
\end{itemize}

\subsection{Adult \& Non-Formal}
\begin{itemize}
    \item Literacy classes
    \item Farmer field schools
    \item Women's groups, cooperatives
\end{itemize}

\subsection{Knowledge \& Information}
\begin{itemize}
    \item Library, reading materials
    \item Bulletin boards, announcements
\end{itemize}

%==============================================================================
\section{Healthcare}
%==============================================================================

\subsection{Primary}
\begin{itemize}
    \item Health posts, sub-health posts
    \item FCHV (Female Community Health Volunteers)
    \item Basic medicines, first aid
    \item Vaccination
\end{itemize}

\subsection{Hospital \& Referral}
\begin{itemize}
    \item District, zonal, regional hospitals
    \item Kathmandu for specialized care
    \item Ambulance, doko, stretcher, helicopter
    \item Lab, X-ray, surgery
\end{itemize}

\subsection{Maternal \& Child}
\begin{itemize}
    \item Antenatal care, delivery (home vs facility)
    \item Birthing centers
    \item Nutrition, immunization
\end{itemize}

\subsection{Specialized}
\begin{itemize}
    \item Dental, eye care (often camp-based)
    \item Mental health (underserved)
    \item Disability, rehabilitation
    \item Elderly, palliative care
\end{itemize}

\subsection{Traditional \& Alternative}
\begin{itemize}
    \item \textit{Dhami/jhankri}, Ayurveda
    \item Herbal medicine
\end{itemize}

\subsection{Pharmacy \& Drugs}
\begin{itemize}
    \item Drug availability, quality
    \item Pharmacy, health post stock
\end{itemize}

%==============================================================================
\section{Administration \& Government}
%==============================================================================

\subsection{Offices}
\begin{itemize}
    \item Village/municipal office (\textit{gaunpalika/nagarpalika})
    \item Ward office (\textit{ward karyalaya})
    \item Ward chairperson, staff
    \item District office access
\end{itemize}

\subsection{Documentation}
\begin{itemize}
    \item Citizenship (\textit{nagarikta})
    \item Birth, death, marriage registration
    \item Land ownership (\textit{lalpurja}), registration
    \item Tax payment
    \item PAN, business registration
    \item Passport, visa
\end{itemize}

\subsection{Justice \& Security}
\begin{itemize}
    \item Police post
    \item Court access
    \item Mediation, dispute resolution
    \item Domestic violence, child protection
\end{itemize}

\subsection{Political \& Civic}
\begin{itemize}
    \item Election booths, voter registration
    \item Public hearings, participation
    \item Citizen assembly, planning
\end{itemize}

\subsection{Programs \& Services}
\begin{itemize}
    \item Census, surveys
    \item Disaster relief
    \item Employment programs
    \item Subsidies: fertilizer, seeds, elderly, disabled, single women
    \item Food security programs
\end{itemize}

%==============================================================================
\section{Communication}
%==============================================================================

\subsection{Traditional \& Basic}
\begin{itemize}
    \item Postal service
    \item Radio (FM, Nepal Radio)
    \item Word of mouth, messengers
    \item Bulletin boards
\end{itemize}

\subsection{Telephony}
\begin{itemize}
    \item Mobile (2G, 3G, 4G)
    \item Landline (rare)
    \item Satellite phone (emergency)
\end{itemize}

\subsection{Internet \& Digital}
\begin{itemize}
    \item Mobile data
    \item Broadband, fiber (rare)
    \item Cyber cafes, community internet
    \item Smartphones, social media
    \item E-governance, online services
    \item Digital payments, mobile banking
\end{itemize}

\subsection{Media}
\begin{itemize}
    \item Television, cable, satellite
    \item Newspapers (where available)
\end{itemize}

%==============================================================================
\section{Commerce, Finance \& Markets}
%==============================================================================

\subsection{Markets}
\begin{itemize}
    \item Local shops, \textit{haat bazaar}
    \item Wholesale, retail
    \item Weekly, seasonal markets
\end{itemize}

\subsection{Finance}
\begin{itemize}
    \item Banks, branches, mobile banking
    \item Cooperatives, savings groups
    \item Microfinance
    \item Credit, loans
    \item Remittance transfer
\end{itemize}

\subsection{Insurance}
\begin{itemize}
    \item Life, health, property
    \item Crop, livestock (emerging)
\end{itemize}

\subsection{Cooperatives}
\begin{itemize}
    \item Agricultural: input supply, marketing
    \item Consumer
    \item Credit
\end{itemize}

%==============================================================================
\section{Sanitation \& Waste}
%==============================================================================

\begin{itemize}
    \item Toilets: pit, pour-flush, septic; open defecation
    \item Waste: composting, burning, dumping
    \item Drainage, flooding
    \item Solid waste management (rare)
\end{itemize}

%==============================================================================
\section{Social, Cultural \& Religious}
%==============================================================================

\begin{itemize}
    \item Temples, \textit{gumbas}, mosques, churches
    \item Community halls (\textit{samaj ghar})
    \item Festivals, \textit{jatras}, ceremonies
    \item Cemeteries, cremation grounds
    \item Sports: football, volleyball, cricket, \textit{dandi biyo}
    \item Youth clubs, women's groups
    \item Music, dance, theater
    \item Cinema (where available)
\end{itemize}

%==============================================================================
\section{Industry, Crafts \& Livelihood}
%==============================================================================

\subsection{Cottage \& Handicraft}
\begin{itemize}
    \item Weaving, knitting
    \item Pottery, metalwork
    \item Bamboo, cane work
    \item Handicraft for market
\end{itemize}

\subsection{Processing \& Manufacturing}
\begin{itemize}
    \item Rice mills, oil mills
    \item Sawmills, brick kilns
    \item Cheese, ghee, jam
    \item Small-scale manufacturing
\end{itemize}

\subsection{Livelihood}
\begin{itemize}
    \item Agriculture, livestock
    \item Wage labor, seasonal migration
    \item Remittance (major)
    \item Business, shopkeeping
    \item Government, NGO employment
    \item Tourism: homestay, trekking, guide
\end{itemize}

%==============================================================================
\section{Environment \& Natural Resources}
%==============================================================================

\begin{itemize}
    \item Forests: community, government; fuel, timber
    \item Grazing land
    \item Soil conservation, landslide prevention
    \item Biodiversity, wildlife
    \item Climate adaptation
    \item Seasonal calendar: monsoon, drought, cold
\end{itemize}

%==============================================================================
\section{Resource Conservation}
%==============================================================================

Conserving food, energy, water, and other resources reduces waste, extends availability, and builds resilience---especially where supply is limited or seasonal.

\subsection{Food Conservation}
\begin{itemize}
    \item \textbf{Cold storage:} Community cold store (electricity or solar-powered) for vegetables, fruits, dairy, meat---extends shelf life, reduces post-harvest loss (20--50\% for horticulture without it; see References), enables off-season supply
    \item \textbf{Evaporative cooling:} \textit{Pot-in-pot}, sand storage---low-tech cooling without electricity
    \item \textbf{Drying:} Sun-drying (\textit{sukeko}), solar dryer---grains, vegetables, fruits, fish, meat (\textit{sukuti}); reduces moisture, prevents spoilage
    \item \textbf{Pickling \& fermentation:} \textit{Achar}, \textit{gundruk}, \textit{sinki}, \textit{kinema}---preserves vegetables, extends use
    \item \textbf{Smoking:} Meat, fish---traditional preservation
    \item \textbf{Salt curing:} Fish, meat
    \item \textbf{Grain storage:} \textit{Bharako} (elevated, rodent-proof), neem leaves, ash---protects from pests, moisture
    \item \textbf{Seed storage:} Cool, dry, sealed---maintains viability
    \item \textbf{Root cellar:} Underground or shaded---potato, onion, ginger
    \item \textbf{Reducing waste:} Proper handling, timely harvest, first-in-first-out
\end{itemize}

\subsection{Energy Conservation}
\begin{itemize}
    \item \textbf{Cooking:} Improved cookstoves (less fuel), pressure cooker, lid on pot, batch cooking
    \item \textbf{Lighting:} LED bulbs, solar lanterns---efficient use of limited power
    \item \textbf{Insulation:} Houses---reduce heating/cooling need
    \item \textbf{Appliance efficiency:} Energy-efficient fridge, fan, pump---where electricity exists
    \item \textbf{Peak avoidance:} Use mills, pumps during off-peak---if grid with load management
    \item \textbf{Biomass:} Sustainable harvest, replanting---don't deplete forest
    \item \textbf{Waste heat:} Biogas slurry for fertilizer; improved chulo heat for water
\end{itemize}

\subsection{Water Conservation}
\begin{itemize}
    \item \textbf{Rainwater harvesting:} Rooftop, tanks---capture monsoon
    \item \textbf{Greywater reuse:} Washing water for garden, livestock
    \item \textbf{Irrigation:} Drip, sprinkler---reduce loss; canal lining
    \item \textbf{Leak repair:} Pipes, taps---prevent waste
    \item \textbf{Spring protection:} Source conservation, recharge
    \item \textbf{Storage:} Tanks, ponds---for dry season
\end{itemize}

\subsection{Material \& Other Resource Conservation}
\begin{itemize}
    \item \textbf{Reuse:} Containers, bags, tools---repair before replace
    \item \textbf{Recycling:} Metal, plastic---where collection exists; composting organic
    \item \textbf{Timber:} Sustainable harvest, efficient use, alternative (bamboo)
    \item \textbf{Soil:} Compost, mulch, terrace---conserve fertility
    \item \textbf{Medicines:} Proper storage, expiry awareness---avoid waste
    \item \textbf{Fuel:} Efficient stoves, no idle burning
\end{itemize}

\subsection{Community-Level Conservation}
\begin{itemize}
    \item \textbf{Grain bank:} Collective storage, buffer for lean season
    \item \textbf{Seed bank:} Preserve local varieties, share
    \item \textbf{Water user group:} Fair allocation, maintenance
    \item \textbf{Forest user group:} Regulated harvest, regeneration
\end{itemize}

%==============================================================================
\section{Security \& Emergency}
%==============================================================================

\begin{itemize}
    \item Natural disasters: earthquake, landslide, flood
    \item Early warning, evacuation
    \item Fire: limited services; community response
    \item Conflict, dispute
    \item Wild animals, snakes
    \item Search \& rescue
    \item Relief distribution
\end{itemize}

%==============================================================================
\section{Vulnerable Groups \& Social Protection}
%==============================================================================

\begin{itemize}
    \item Elderly care, pension
    \item Orphan, child protection
    \item Disabled support
    \item Widows, single women
    \item Dalit, marginalized
    \item Indigenous, ethnic
    \item Refugees (if applicable)
\end{itemize}

%==============================================================================
\section{Technology \& Repair}
%==============================================================================

\begin{itemize}
    \item Mobile repair
    \item Battery charging
    \item Solar installation, repair
    \item Agricultural machinery repair
    \item Bicycle, motorcycle repair
\end{itemize}

%==============================================================================
\section{Public Infrastructure}
%==============================================================================

\begin{itemize}
    \item Roads, bridges
    \item Street lighting
    \item Public spaces, playgrounds
    \item Warehouses, storage
    \item Parking, bus stops
\end{itemize}

%==============================================================================
\section{Legal \& Rights}
%==============================================================================

\begin{itemize}
    \item Marriage, divorce
    \item Inheritance, property
    \item Labor rights
    \item Child rights
    \item Women's rights
    \item Land rights, tenancy
\end{itemize}

%==============================================================================
\section{Identity \& Belonging}
%==============================================================================

\begin{itemize}
    \item Caste, ethnicity
    \item Language, mother tongue
    \item Religion
    \item Voter ID
    \item Community membership
\end{itemize}

%==============================================================================
\section{100\% Fully Independent Village: Complete Self-Sufficiency Checklist}
%==============================================================================

What must a village have \textit{within its boundaries} to be 100\% fully independent---requiring no external inputs, services, or markets for survival and functioning? This section catalogs every need for maximum self-reliance. Some items are aspirational; some have \textbf{hard dependencies} (noted) that may require minimal external trade or are nearly impossible to produce locally.

\subsection{Water Independence}
\begin{itemize}
    \item Own water source: spring, well, river, or sufficient rainwater harvesting
    \item Own gravity-fed or locally pumped distribution system
    \item Own purification: boiling, sand filter, solar disinfection---no external chemicals
    \item Own irrigation: \textit{kulo}, ponds, lift from local stream
    \item Own maintenance capacity: masons, plumbers
    \item Water user committee for governance
\end{itemize}

\subsection{Energy Independence}
\begin{itemize}
    \item \textbf{Cooking:} Firewood from community forest (sustainably managed), biogas from livestock, agricultural residue---no LPG, kerosene
    \item \textbf{Electricity:} Micro-hydro, pico-hydro, or solar---no grid dependency
    \item \textbf{Mechanical:} Water mills (\textit{ghatta}), animal power, human power---or electric mills from local generation
    \item \textbf{Lighting:} Solar, fire, oil from local seeds (mustard, etc.)
    \item \textbf{Hard dependency:} Solar panels, batteries, micro-hydro turbines---typically imported; for true independence would need local manufacture or long-life equipment
\end{itemize}

\subsection{Food Independence}
\begin{itemize}
    \item \textbf{Seeds:} Own seed saving, community seed bank---no hybrid dependency
    \item \textbf{Fertilizer:} Compost, manure, green manure, ash---no chemical fertilizer
    \item \textbf{Pesticides:} Neem, local botanicals, cultural methods---no external chemicals
    \item \textbf{Crops:} Full diversity: staples (rice/maize/millet), pulses, vegetables, fruits, spices---sufficient for year-round diet
    \item \textbf{Livestock:} Buffalo/cows (milk, manure, draft), goats, chickens, pigs---for protein, eggs, manure
    \item \textbf{Storage:} \textit{Bharako}, traditional methods---no cold chain
    \item \textbf{Processing:} Own rice mill, oil press, flour mill---powered locally
    \item \textbf{Preservation:} Drying, pickling, smoking, fermentation
    \item \textbf{Hard dependency:} Salt (unless local salt source, e.g., Himalayan); iodized salt for health
\end{itemize}

\subsection{Shelter Independence}
\begin{itemize}
    \item Own building materials: stone, mud, wood, bamboo, thatch, slate
    \item Own masons, carpenters, thatchers
    \item Own repair and construction capacity
    \item Earthquake-resistant techniques (local knowledge)
\end{itemize}

\subsection{Clothing \& Textile Independence}
\begin{itemize}
    \item Own fiber: cotton, wool, hemp, nettle, silk (if sericulture)
    \item Own spinning, weaving
    \item Own tailoring, repair
    \item Own dye: natural dyes from plants, minerals
    \item Own footwear: leather (if livestock), rubber-soled (hard dependency)
\end{itemize}

\subsection{Transport Independence (Internal)}
\begin{itemize}
    \item Foot trails, bridges (locally built)
    \item Pack animals: mules, horses, yaks---bred and maintained locally
    \item Porters (human labor)
    \item No need for motor vehicles if no external travel
    \item Boats/rafts if river within village
\end{itemize}

\subsection{Education Independence}
\begin{itemize}
    \item Own primary school, teachers from village (trained or apprenticed)
    \item Own secondary school (or sufficient population to justify)
    \item Own vocational training: apprenticeship within village
    \item Own library, books (printed---hard dependency; or oral tradition, handwritten)
    \item Curriculum: state-approved or community-defined
    \item Own school management, no external funding
\end{itemize}

\subsection{Healthcare Independence}
\begin{itemize}
    \item Own health worker: trained FCHV, midwife, or community health worker
    \item Own pharmacy: herbal medicine, Ayurveda, basic allopathic (stocked---hard dependency)
    \item Own birthing capacity: skilled birth attendant
    \item Own capacity for: wounds, fractures, common illnesses, dental extraction
    \item Own mental health: community support, traditional healers (\textit{dhami/jhankri})
    \item Own vaccination: cold chain impossible without electricity; or accept disease risk
    \item \textbf{Hard dependency:} Antibiotics, vaccines, surgical equipment, complex diagnostics---cannot be produced locally; 100\% independence means accepting limits or minimal strategic stockpile from one-time trade
\end{itemize}

\subsection{Governance \& Administration Independence}
\begin{itemize}
    \item Own village council, ward committee---elected or customary
    \item Own dispute resolution: mediation, \textit{panchayat}-style
    \item Own record-keeping: births, deaths, land---written or oral
    \item Own enforcement: community norms, local guards
    \item \textbf{State interface:} Citizenship, national ID, land registration---minimal; village can maintain internal records; full independence implies village as de facto autonomous unit (historical \textit{rajya}, \textit{thum})
\end{itemize}

\subsection{Justice \& Security Independence}
\begin{itemize}
    \item Own dispute resolution for civil, family, property
    \item Own security: no external police; community vigilance
    \item Own disaster response: search, rescue, relief from local stocks
    \item Own fire response: bucket brigades, community
\end{itemize}

\subsection{Communication Independence}
\begin{itemize}
    \item Internal: word of mouth, village crier, notice boards
    \item Own postal runner to neighboring villages (if any exchange needed)
    \item Own radio transmitter (if electricity and equipment---hard dependency)
    \item No mobile, internet---or own community network (very advanced)
\end{itemize}

\subsection{Commerce \& Finance Independence}
\begin{itemize}
    \item Barter, labor exchange (\textit{parma}, \textit{parma})
    \item Own grain bank, seed bank
    \item Own savings group, rotating fund
    \item No external currency dependency for internal transactions
    \item No bank, no remittance---economy is closed
\end{itemize}

\subsection{Sanitation \& Waste Independence}
\begin{itemize}
    \item Own toilets: pit, compost---no external septic services
    \item Own waste: composting, reuse, minimal dumping
    \item Own drainage: constructed and maintained locally
\end{itemize}

\subsection{Industry, Craft \& Repair Independence}
\begin{itemize}
    \item \textbf{Blacksmith:} Tools, implements, repair---ore (hard: need local iron or recycling)
    \item \textbf{Carpenter:} Furniture, construction, repair
    \item \textbf{Potter:} Utensils, storage
    \item \textbf{Weaver, tailor:} As above
    \item \textbf{Miller:} Rice, flour, oil
    \item \textbf{Bicycle/motorcycle repair:} Only if vehicles used; parts (hard dependency)
    \item \textbf{Solar/electrical repair:} If solar used; components (hard dependency)
\end{itemize}

\subsection{Raw Materials \& Hard Dependencies}
\begin{itemize}
    \item \textbf{Salt:} Himalayan rock salt in some areas; otherwise trade
    \item \textbf{Iron/steel:} Mining rare; recycling from scrap; otherwise trade
    \item \textbf{Soap:} Can make from ash + local oil
    \item \textbf{Paper:} Can make from lokta, recycled; or palm leaf, bark
    \item \textbf{Glass:} Cannot produce locally; use oiled paper, open windows
    \item \textbf{Rubber, plastic:} Cannot produce; minimize use
    \item \textbf{Medicines:} Herbal for most; critical allopathic---strategic reserve or accept risk
\end{itemize}

\subsection{Population \& Demographics}
\begin{itemize}
    \item Sufficient population to fill all roles: farmers, blacksmith, carpenter, teacher, health worker, etc.
    \item Marriage pool: internal or arranged with neighboring villages (reduces strict independence)
    \item Birth rate to sustain population
    \item Care for elderly, disabled, orphans---within community
\end{itemize}

\subsection{Social, Cultural \& Religious}
\begin{itemize}
    \item Own temples, \textit{gumbas}, places of worship
    \item Own community hall
    \item Own festivals, rituals, ceremonies
    \item Own entertainment: music, dance, storytelling
    \item Own cremation/burial grounds
\end{itemize}

\subsection{Environment \& Sustainability}
\begin{itemize}
    \item Community forest: sustainably managed for fuel, timber
    \item Grazing land for livestock
    \item Soil conservation, terrace maintenance
    \item No over-extraction: water, forest, soil
\end{itemize}

\subsection{Time, Calendar \& Measurement}
\begin{itemize}
    \item \textbf{Timekeeping:} Sundial, water clock, stars, rooster---for work rhythm
    \item \textbf{Calendar:} Lunar (\textit{tithi}), solar---for planting, festivals, rituals
    \item \textbf{Agricultural calendar:} Who knows when to plant? Weather signs, traditional knowledge
    \item \textbf{Measurement:} \textit{Pathi}, \textit{mana}, \textit{dharni} for grain; \textit{haat}, \textit{kosh} for length---standardization within village
    \item \textbf{Value:} How is labor valued in \textit{parma}? How much grain for a day's work? Custom, negotiation
\end{itemize}

\subsection{Surplus, Buffer \& Risk}
\begin{itemize}
    \item \textbf{Famine reserve:} How many months of grain? One bad year? Two? Strategic storage
    \item \textbf{Winter buffer:} High mountains---6+ months snow; fuel, food, fodder stored
    \item \textbf{Disaster stock:} Earthquake, flood---immediate relief from local stores
    \item \textbf{Crop failure:} Hail, drought, pest---diversity, fallback crops
    \item \textbf{Livestock epidemic:} Quarantine, culling, herbal treatment---no vet
\end{itemize}

\subsection{Geographic \& Resource Reality}
\begin{itemize}
    \item \textbf{All within boundary:} Arable land, forest, water source, pasture, settlement---must all exist within village territory
    \item \textbf{Altitude:} Hill villages often need land at multiple elevations (rice low, potato high); scattered plots
    \item \textbf{Transhumance:} Seasonal movement of livestock---grazing rights
    \item \textbf{Water source:} Spring, river---must be within or legally secured; upstream villages can cut supply
    \item \textbf{Quarries:} Stone, clay, lime---local deposits or trade
\end{itemize}

\subsection{Skill Redundancy \& Succession}
\begin{itemize}
    \item \textbf{No single point of failure:} If blacksmith dies, who replaces? Multiple apprentices
    \item \textbf{Knowledge transfer:} Oral, written, apprenticeship---before elder dies
    \item \textbf{Leadership succession:} How is next leader chosen? Hereditary, election, rotation
    \item \textbf{Chicken-egg:} Blacksmith needs anvil; potter needs wheel---initial tools from somewhere; recycling, inheritance
\end{itemize}

\subsection{Death, Ritual \& Continuity}
\begin{itemize}
    \item \textbf{Cremation/burial:} Wood for pyre, land for burial---who provides?
    \item \textbf{Ritual:} Who performs? Priest, lama---full-time or part-time; how compensated?
    \item \textbf{Inheritance:} Land, house, tools---rules for division; fragmentation over generations
    \item \textbf{Orphan:} Who raises? Kinship, adoption
    \item \textbf{Widow:} Property rights, remarriage---custom
\end{itemize}

\subsection{Conflict, Crime \& Punishment}
\begin{itemize}
    \item \textbf{Dispute:} Civil, family, property---mediation; what if parties refuse?
    \item \textbf{Crime:} Murder, theft, rape---what punishment? Fine, exile, corporal? No jail without guards, food
    \item \textbf{Factionalism:} Family feud, caste conflict---can split village
    \item \textbf{Boundary dispute:} With neighboring village---grazing, forest, water; no external arbiter
\end{itemize}

\subsection{Defense \& External Threat}
\begin{itemize}
    \item \textbf{Bandits, raiders:} Historical reality---walls, vigilance, weapons?
    \item \textbf{Wild animals:} Leopard, bear, wild boar---crops, livestock, humans; how to deter?
    \item \textbf{Neighboring village:} Conflict over resources---defense or negotiation
\end{itemize}

\subsection{Healthcare: The Limits}
\begin{itemize}
    \item \textbf{Snakebite:} Antivenom---cannot produce; herbal alternatives limited
    \item \textbf{Complicated childbirth:} Breech, hemorrhage---traditional skills can save some; many die
    \item \textbf{Appendicitis, trauma:} Surgery---no surgeon; accept death
    \item \textbf{Cataract, severe dental:} Blindness, chronic pain---camp-based care requires external visit
    \item \textbf{Mental illness:} Severe---restraint, care; no psychiatric drugs
    \item \textbf{Contagious epidemic:} Quarantine, isolation---who enforces?
\end{itemize}

\subsection{Population \& Demographics: The Numbers}
\begin{itemize}
    \item \textbf{Minimum viable:} Estimate 200--500+ to fill all roles with redundancy
    \item \textbf{Genetic diversity:} Inbreeding if too closed; marriage with neighbors reduces independence
    \item \textbf{Birth/death balance:} Sustain population; family planning---natural methods
    \item \textbf{Disability, chronic illness:} Care within family/community---productivity loss
\end{itemize}

\subsection{Social Obligation \& Reciprocity}
\begin{itemize}
    \item \textbf{Hospitality:} Cultural duty to feed guests---uses surplus
    \item \textbf{Festival feasts:} Who provides? Rotating, pooling
    \item \textbf{Gift economy:} Marriage, death, festivals---reciprocity, social capital
    \item \textbf{Common labor:} \textit{Shramdan}---trail, canal, temple---who organizes? Who participates?
\end{itemize}

\subsection{Division of Labor: Caste, Gender}
\begin{itemize}
    \item \textbf{Traditional:} Blacksmith (\textit{kami}), tailor (\textit{damai}), etc.---caste-based; does independence require breaking or retaining?
    \item \textbf{Gender:} Who fetches water? Who plows? Who weaves? Who cares for children?
    \item \textbf{Child labor:} At what age do children work? School vs. field
\end{itemize}

\subsection{Childcare, Elderly \& Disability}
\begin{itemize}
    \item \textbf{Childcare:} When parents in field---grandmother, older sibling, communal
    \item \textbf{Elderly:} Bedridden---who feeds, bathes? Family, rotation
    \item \textbf{Disabled:} Born or acquired---who supports? No pension; family, charity
\end{itemize}

\subsection{Storage \& Preservation}
\begin{itemize}
    \item \textbf{Grain:} Rats, weevils---how to protect? \textit{Bharako} design, neem leaves, smoke---studies report 5--20\% cereal storage loss in Nepal; hermetic storage adoption remains low (see References)
    \item \textbf{Seed viability:} Saving year to year---maintain purity, germination
\end{itemize}

\subsection{Pollution \& Waste}
\begin{itemize}
    \item \textbf{Blacksmith:} Smoke, slag---where?
    \item \textbf{Tannery:} If leather---waste, smell
    \item \textbf{Latrine:} Pit filling---where next? Rotation
\end{itemize}

\subsection{Beauty, Meaning \& Spirit}
\begin{itemize}
    \item \textbf{Art:} Mural, carving, craft---for meaning, not just function
    \item \textbf{Entertainment:} Music, dance, storytelling---who performs? When?
    \item \textbf{Instruments:} Drum, flute---who makes? Who plays?
    \item \textbf{Why:} Purpose beyond survival---ritual, festival, belonging
\end{itemize}

\subsection{The Hard Truths}
\begin{enumerate}
    \item \textbf{100\% is impossible} without accepting death, disability, and suffering that modern medicine could prevent.
    \item \textbf{Geographic luck} matters: a village without local iron, salt, or sufficient water cannot be independent.
    \item \textbf{Population} must be large enough for all skills plus redundancy---small villages cannot.
    \item \textbf{Marriage} with neighbors breaks strict independence; inbreeding breaks health.
    \item \textbf{Surplus} is essential---for bad years, festivals, guests---but surplus requires good land, labor, and luck.
    \item \textbf{Conflict} without external arbiter can destroy a village from within.
    \item \textbf{Independence} was historically possible when villages were \textit{rajya}, \textit{thum}---semi-autonomous. The modern state, market, and aspiration have made it rare.
\end{enumerate}

\subsection{Summary: Independence vs.\ Reality}
A 100\% fully independent village is theoretically possible only if: (1) it produces or harvests everything from local resources, (2) it has sufficient population and skill diversity, (3) it accepts \textbf{hard dependencies} for a few items (salt, iron, critical medicines) via minimal trade, or (4) it accepts higher risk (disease, crop failure) without external buffer. In practice, \textit{maximum self-reliance} is the achievable goal---minimizing external dependency while maintaining dignity and resilience.

%==============================================================================
\section{Summary: The Continuum}
%==============================================================================

\begin{description}
    \item[Fully Remote] Local water, firewood, subsistence agriculture, foot/mule transport, health posts, primary schools, ward office, radio, occasional mobile. No grid, no roads, no internet.
    \item[Partially Connected] Micro-hydro or solar, trail/ropeway, secondary school access, district hospital referral, weekly market, patchy mobile.
    \item[Fully Connected] Roads, grid, internet, banks, higher education, better healthcare, LPG, vehicles, e-services.
    \item[100\% Independent] Maximum self-reliance: local water, fuel, food, shelter, education, healthcare, governance, industry---minimal external trade for salt, iron, critical medicines.
\end{description}

Understanding this continuum helps in planning interventions, policy, and development projects tailored to each village's reality.

%==============================================================================
\section{References and Further Reading}
%==============================================================================

The following research articles, books, and reports inform this framework. Links and citations are provided for deeper study.

\subsection{Infrastructure \& Wellbeing}
\begin{itemize}
    \item \textbf{Access to infrastructure and human well-being: evidence from rural Nepal.} Examines how roads, water, irrigation, electricity, and other infrastructure affect wellbeing; finds roads highest priority, followed by drinking water and irrigation; remote areas show 50--65\% lower per capita income, higher malnutrition and dropout. \url{https://scispace.com/pdf/access-to-infrastructure-and-human-well-being-evidence-from-4uuy9g9dkt.pdf}
    \item \textbf{Rural Living in Nepal: Housing, Water \& Waste in Paanchkhal.} NPRC Journal of Multidisciplinary Research. Housing (Pakka/Kachcha), water sources, purification, latrines, waste disposal. \url{https://nepjol.info/index.php/nprcjmr/article/view/78324}
\end{itemize}

\subsection{Water \& Irrigation}
\begin{itemize}
    \item \textbf{Does Rural Water System Design Matter? A Study of Productive Use of Water in Rural Nepal.} Water 11(10), 1978. Productive use of water beyond drinking. \url{https://www.mdpi.com/2073-4441/11/10/1978}
    \item \textbf{Functionality of Rural Community Water Supply Systems and Collective Action: Guras Rural Municipality, Karnali.} \url{https://nppr.org.np/index.php/journal/article/view/26/53}
    \item \textbf{Sulitar Irrigation System: Struggle from Poverty to Prosperity.} Hydro Nepal. Farmer-managed irrigation, participatory approach, community ownership. \url{https://www.nepjol.info/index.php/HN/article/view/13276}
    \item \textbf{Farmers in Nepal reap benefits of rehabilitated farmers-managed irrigation systems.} Asian Development Bank. FMIS rehabilitation, crop intensification, income gains. \url{https://www.adb.org/results/farmers-nepal-reap-benefits-rehabilitated-farmers-managed-irrigation-systems}
    \item \textbf{Improving irrigation governance and management in Nepal.} Virginia Tech. 20+ years of research on community resource management. \url{https://vtechworks.lib.vt.edu/items/cbc1aa78-7c13-4404-bfbf-2552edcc72e7}
    \item \textbf{Irrigation in Nepal: Opportunities and Constraints.} Third World Centre. \url{https://thirdworldcentre.org/wp-content/uploads/2020/07/RPP-Dec-1-1989-Irrigation-in-Nepal-Opportunities-and-Constraints.pdf}
\end{itemize}

\subsection{Food Security \& Post-Harvest}
\begin{itemize}
    \item \textbf{Karki TB, Sah SK, Thapa RB, McDonald AJ, Davis AS (2015) Identifying Pathways for Improving Household Food Self-Sufficiency Outcomes in the Hills of Nepal.} PLoS ONE 10(6): e0127513. CART analysis; agronomic management, maize-fingermillet yields, improved cultivars as primary drivers. \url{https://doi.org/10.1371/journal.pone.0127513}
    \item \textbf{Post-harvest practices of horticultural crops in Nepal: Issues and management.} 20--50\% post-harvest loss; cold storage, packaging. \url{https://www.academia.edu/66093653/Post_harvest_practices_of_horticultural_crops_in_Nepal_Issues_and_management}
    \item \textbf{On-Farm Grain Storage and Challenges in Bagmati Province, Nepal.} Sustainability. 5--20\% storage loss; pests, monsoon; hermetic storage adoption low. \url{https://www.cgiar.org/research/publication/on-farm-grain-storage-and-challenges-in-bagmati-province-nepal/}
    \item \textbf{Assessing the Impact of Hermetic Storage Technology on Post-harvest Losses among Maize Farmers in Nepal.} Preprints. \url{https://www.preprints.org/manuscript/202411.1421/v1}
\end{itemize}

\subsection{Energy}
\begin{itemize}
    \item \textbf{Techno-Economic Modelling of Micro-Hydropower Mini-Grids in Nepal.} Energies. Financial sustainability, electric cooking, productive use. \url{https://www.mdpi.com/2071-1050/14/16/4232}
    \item \textbf{Identifying the best decentralized renewable energy system for rural electrification in Nepal.} Journal of Asian Rural Studies. Micro-hydro ranked best, then solar. \url{http://pasca.unhas.ac.id/ojs/index.php/jars/article/view/2097}
    \item \textbf{Performance evaluation of solar PV mini grid system in Nepal: Thabang and Sugarkhal.} Frontiers in Energy Research. \url{https://www.frontiersin.org/journals/energy-research/articles/10.3389/fenrg.2023.1321945/full}
    \item \textbf{Sustainability assessment of a micro hydropower plant in Nepal.} Energy, Sustainability and Society. \url{https://energsustainsoc.biomedcentral.com/articles/10.1186/s13705-018-0147-2}
\end{itemize}

\subsection{Health \& Education}
\begin{itemize}
    \item \textbf{Perception of Health and Care Seeking Behaviors in High-Altitude Villages of Rural Nepal.} HPROSPECT. Traditional healers, cultural beliefs, accessibility. \url{https://nepjol.info/index.php/HPROSPECT/article/view/24103}
    \item \textbf{Assessing the health impacts of a Comprehensive Rural Health Project in Nepal.} PLOS Global Public Health. Village Alive Project, Dalit villages. \url{https://journals.plos.org/globalpublichealth/article?id=10.1371/journal.pgph.0004458}
\end{itemize}

\subsection{Rural Development \& Transport}
\begin{itemize}
    \item \textbf{Rural Transportation Infrastructure in Low- and Middle-Income Countries.} Sustainability. Impacts, interventions, links to agriculture and health. \url{https://www.mdpi.com/2071-1050/14/4/2149}
    \item \textbf{Handbook of Rural Development.} Comprehensive rural development scope.
    \item \textbf{Infrastructure for Smart Villages.} Routledge. Fit-for-purpose rural infrastructure, SDG alignment.
    \item \textbf{UN Handbook on Infrastructure Asset Management.} Managing infrastructure for sustainable development.
\end{itemize}

\end{document}
